\centerline{\bf \Huge Оценочная теория чисел}

\begin{flushright}
        \scriptsize{Эпизод 9. Действовать синхронно}
\end{flushright}

\problem{1}
Докажите, что $(a b+1, b c+1, c a+1) \leqslant \frac{a+b+c}{3}$ для попарно различных натуральных $a, b, c$.

\problem{2}
Найдите все натуральные числа $a, b$ и $c$ такие, что $\frac{(a b-1)(a c-1)}{b c}$ является квадратом целого числа.

\problem{3}
Найдите все пары натуральных чисел ($n, p$ ), где $p$ простое и

$$
1+2+\ldots+n=3 \cdot\left(1^2+2^2+\ldots+p^2\right) .
$$


\problem{4}
Можно ли выписать 90 различных трёхзначных чисел в ряд в порядке возрастания так, чтобы произведение любых двух подряд идущих чисел делилось на предыдущее число?

\problem{5}
Натуральные числа $a, b, c, d, e$ и $f<a$ таковы, что $a b d+1$ делится на $c, a c e+1$ делится на $b, b c f+1$ делится на $a$. Докажите, что если $d / c<1-e / b$, то $d / c<1-f / a$.

\problem{6}
Найдите все тройки натуральных чисел $(a, b, c)$ такие, что $a^3+b^3+c^3=(a b c)^2$.

\problem{7}
Дано простое число $p>3$. Натуральные числа $a, b, c$ и $d$ таковы, что

$$
a+b=c+d=p,
$$


а число $a b c d$ - точный квадрат. Докажите, что какие-то два из чисел $a, b, c$ и $d$ совпадают.